\documentclass[12pt, a4paper, oneside]{ctexart}
\usepackage{amsmath, amsthm, amssymb, bm, xcolor, tikz, framed, graphicx, hyperref, mathrsfs}

\title{Soild State Physics Homework (Week 7-8)}
\author{Tang Yi \\ [学号已脱敏]}
\usepackage[a4paper, left=2cm, right=2cm, top=2cm, bottom=2cm]{geometry}
\linespread{1.3} % 行距更紧凑
\definecolor{shadecolor}{RGB}{241, 241, 255}

\newcounter{problemname}
\newenvironment{problem}
  {\begin{shaded}\stepcounter{problemname}\par\noindent\textbf{题目 \theproblemname. }}
  {\end{shaded}\par}
\newenvironment{solution}
  {\par\noindent\textbf{题目 \theproblemname 的解答. }}
  {\par}
\newenvironment{note}
  {\par\noindent\textbf{题目 \theproblemname 的注记. }}
  {\par}

\begin{document}
\maketitle
\begin{problem}
已知 NaCl 晶体平均每对离子的相互作用能为
\[
U(r) = -\alpha q^2 / r + \beta / r^n
\]
其中马德隆常数 $\alpha = 1.75$，$n = 9$，平衡离子间距 $r_0 = 2.82\,\text{\AA}$。
\begin{enumerate}
\item 试求离子在平衡位置附近的振动频率。
\item 计算与该频率相当的电磁波的波长，并与 NaCl 红外吸收频率的测量值 $61\,\mu$m 进行比较。
\end{enumerate}
\end{problem}

\begin{solution}
(1) 在平衡位置附近作小振动展开：
\[
\omega = \sqrt{\frac{1}{\mu}\left.\frac{d^2 U}{dr^2}\right|_{r_0}}
\]
先求导：
\[
\frac{dU}{dr} = \alpha q^2 \frac{1}{r^2} - n\beta \frac{1}{r^{n+1}}
\]
平衡条件：
\[
\frac{dU}{dr}\bigg|_{r_0}=0 \Rightarrow \beta = \frac{\alpha q^2}{n} r_0^{n-1}
\]
二阶导数：
\[
\frac{d^2U}{dr^2} = -2\alpha q^2 \frac{1}{r^3} + n(n+1)\beta \frac{1}{r^{n+2}}
\]
代入 $\beta$：
\[
\left.\frac{d^2U}{dr^2}\right|_{r_0}
= \frac{\alpha q^2}{r_0^3}(n-1)
\]
故：
\[
\omega = \sqrt{\frac{\alpha q^2 (n-1)}{\mu r_0^3}}
\]
计算得：

\[
\omega \approx 3.5 \times 10^{13}\ \text{s}^{-1}
\]
(2) 电磁波波长：
\[
\lambda = \frac{2\pi c}{\omega}
\]
得到：
\[
\lambda \approx 5.4\times 10^{-5}\,\text{m} = 54\,\mu\text{m}
\]
代入典型数值（$\mu$ 为 Na 和 Cl 的约化质量）可得数量级在几十微米，与 $61\,\mu$m 同一量级，说明模型合理。
\end{solution}

\begin{problem}
求出一维单原子链的频率分布函数 $\rho(\omega)$。
\end{problem}

\begin{solution}
色散关系：
\[
\omega = 2\sqrt{\frac{K}{M}} \left|\sin\frac{qa}{2}\right|
\]
态密度定义：
\[
\rho(\omega) = \frac{L}{2\pi}\left|\frac{dq}{d\omega}\right|
\]
求导：
\[
\frac{d\omega}{dq} = \sqrt{\frac{K}{M}} a \cos\frac{qa}{2}
\]
因此：
\[
\rho(\omega) = \frac{L}{\pi a}\frac{1}{\sqrt{\omega_m^2 - \omega^2}}
\]
其中 $\omega_m = 2\sqrt{K/M}$。
\end{solution}

\begin{problem}
设三维晶格的光学振动在 $q=0$ 附近的长波极限有
\[
\omega(q) = \omega_0 - A q^2
\]
求证：频率分布函数为
\[
f(\omega) = \frac{V}{4\pi^2} \frac{1}{A^{3/2}} (\omega_0 - \omega)^{1/2}, \quad \omega < \omega_0,
\]
\[
= 0, \quad \omega > \omega_0.
\]
\end{problem}

\begin{solution}
三维态密度：
\[
f(\omega) = \frac{V}{(2\pi)^3} \int dS_q \frac{1}{|\nabla_q \omega|}
\]
等频面为球面：
\[
q = \sqrt{\frac{\omega_0 - \omega}{A}}
\]
球面积：
\[
4\pi q^2
\]
梯度：
\[
|\nabla_q \omega| = 2Aq
\]
因此：
\[
f(\omega) = \frac{V}{(2\pi)^3} \cdot \frac{4\pi q^2}{2Aq}
= \frac{V}{4\pi^2 A^{3/2}}(\omega_0 - \omega)^{1/2}
\]
且当 $\omega>\omega_0$ 无解，因此为 0。
\end{solution}

\begin{problem}
有 $N$ 个相同原子组成面积为 $S$ 的二维晶格，在德拜近似下计算比热，并论述在低温极限比热正比于 $T^2$。
\end{problem}

\begin{solution}
二维态密度：
\[
g(\omega) \propto \omega
\]
内能：
\[
U = \int_0^{\omega_D} \frac{\hbar \omega}{e^{\hbar\omega/k_BT}-1} g(\omega)\,d\omega
\]
低温近似：
\[
U \propto T^3
\]
因此比热：
\[
C = \frac{dU}{dT} \propto T^2
\]
\end{solution}

\begin{problem}
写出量子谐振子系统的自由能，证明在经典极限，自由能为
\[
F \simeq U_0 + k_B T \sum_q \ln\left(\frac{\hbar \omega_q}{k_B T}\right).
\]
\end{problem}

\begin{solution}
单模配分函数：
\[
Z = \frac{e^{-\beta \hbar\omega/2}}{1 - e^{-\beta \hbar\omega}}
\]
自由能：
\[
F = -k_BT \ln Z
= \frac{1}{2}\hbar\omega + k_BT \ln(1 - e^{-\beta\hbar\omega})
\]
高温极限：
\[
1 - e^{-x} \approx x
\]
得到：
\[
F \approx \frac{1}{2}\hbar\omega + k_BT \ln(\beta\hbar\omega)
\]
对所有模式求和：
\[
F \simeq U_0 + k_B T \sum_q \ln\left(\frac{\hbar \omega_q}{k_B T}\right)
\]
\end{solution}

\begin{problem}
设晶体中每个振子的零点振动能为 $\frac{1}{2}\hbar \omega$，试用德拜模型求晶体的零点振动能。
\end{problem}

\begin{solution}
总零点能：
\[
E_0 = \sum_q \frac{1}{2}\hbar\omega
\]
转为积分：
\[
E_0 = \int_0^{\omega_D} \frac{1}{2}\hbar\omega g(\omega)\,d\omega
\]
三维德拜模型：
\[
g(\omega) \propto \omega^2
\]
因此：
\[
E_0 \propto \int_0^{\omega_D} \omega^3 d\omega \propto \omega_D^4
\]
最终结果：
\[
E_0 = \frac{9}{8} N \hbar \omega_D
\]
\end{solution}

\end{document}